% ============================================================================
% Fichier  : partie2/P02_C07_chronologie_correlation.tex
% Titre    : Chronologie et Corrélation
% Auteur   : MINKA MI NGUIDJOÏ Thierry Emmanuel
% Version  : 1.0 -- novembre 2024
% Corrigés : corriges/Corr_C07.tex
% ============================================================================
% Sources : MITRE ATT&CK (tactiques) ; chaîne d'intrusion de Lockheed Martin
% (2011) ; outils de super-chronologie (log2timeline/Plaso, présentés sans
% citation).
% Faits nouveaux (bible) : temps de présence 11 j 7 h 42 min (02/03 23 h 48
% -> 14/03 7 h 30) ; trou du 3 au 8 mars ; dérive SRV-AD01 +2 s, passerelle
% -47 s (horloge non synchronisée).
% ============================================================================

\chapter{Chronologie et Corrélation}
\label{chap:chronologie}
\chaptermark{Chronologie et corrélation}

\epigraph{%
    « Une chronologie n'est pas une liste d'événements.\\
    C'est une liste d'événements dont on a vérifié les horloges. »%
}{Principe du cours}

\niveauChapitre{Intermédiaire}{%
    Maîtrise des Chapitres~\ref{chap:disque} à~\ref{chap:reseau}
}

\begin{boxobjectifs}
\begin{itemize}
    \item Construire une chronologie unifiée à partir de sources aux
          horloges et aux formats différents.
    \item Normaliser les heures et corriger une dérive mesurée.
    \item Corréler des événements par des pivots~: adresse, compte,
          empreinte, heure.
    \item Repérer les trous d'une chronologie et les coter.
    \item Calculer un temps de présence et le rattacher aux tactiques
          d'une attaque.
\end{itemize}
\end{boxobjectifs}

\bigskip

% ============================================================================
\section*{Avant de commencer : quarante-sept secondes}
% ============================================================================

Yaoundé, 26 mars 2025. M.~FOTSO assemble les événements établis par le
disque, la mémoire et le réseau. Une incohérence apparaît aussitôt. Selon le
journal de la passerelle, la session de \texttt{svc\_backup} qui précède
l'installation du service sur SRV-AD01, le 9~mars, s'ouvre à 1~h~35~min~41~s.
Selon le journal système de SRV-AD01, le service est créé à
1~h~35~min~09~s. L'attaquant aurait installé son outil trente-deux secondes
avant d'ouvrir sa session.

Mme~MBARGA ne conclut pas à une falsification. Elle demande la dérive de
chaque horloge. Les serveurs du siège se synchronisaient sur le serveur de
temps interne~; la passerelle, installée par un prestataire, ne le faisait
pas. Mesurée le 14~mars, son horloge avançait de 47~secondes.

\begin{boxalert}
\textbf{Violation typique.} Fusionner les heures de plusieurs sources sans
vérifier leurs fuseaux ni leurs dérives, puis tirer des conclusions de
l'ordre obtenu.

\smallskip
\textbf{Règle~:} chaque source reçoit, avant toute fusion, son fuseau, sa
dérive mesurée et sa précision. La chronologie se construit en UTC
corrigé~; un ordre entre deux événements plus proches que l'incertitude des
horloges n'est pas établi.
\end{boxalert}

\separateur

% ============================================================================
\section{Normaliser les Heures}
\label{sec:c07-normaliser}
% ============================================================================

\begin{tableaucours}{L{3.6cm}cccL{2.6cm}}{%
    \tableauHeader{Source} &
    \tableauHeader{Fuseau} &
    \tableauHeader{Dérive} &
    \tableauHeader{Précision} &
    \tableauHeader{Synchronisée}
}
Passerelle d'accès & UTC+1 & $+47$~s & 1~s & Non \\
SRV-AD01 (journaux) & UTC+1 & $+2$~s & 1~s & Oui \\
SRV-CPT01 (journaux, NTFS) & UTC+1, UTC & $+2$~s & 1~s, 100~ns & Oui \\
Sonde réseau & UTC & $0$~s & 1~µs & Oui (source publique) \\
Interface de paiement & UTC & Inconnue & 1~s & \NEK{} \\
\end{tableaucours}
\vspace{-0.8em}
\begin{center}{\small\itshape Tableau~7.1. Les horloges de l'affaire NYANGA}\end{center}

La dérive se mesure au moment de l'acquisition, en comparant l'heure de la
machine à une référence, et se suppose constante sur la période, ce qui est
une hypothèse à écrire. Corrigée de 47~secondes, l'ouverture de session sur
la passerelle passe à 1~h~34~min~54~s, quinze secondes avant la création du
service~: l'ordre redevient cohérent.

\begin{boxinfo}
\textbf{Préparer une source avant fusion}
\begin{enumerate}
    \item Identifier le fuseau dans lequel la source écrit ses heures.
    \item Mesurer ou rechercher la dérive de son horloge.
    \item Noter sa précision (seconde, milliseconde).
    \item Convertir en UTC et corriger la dérive, en gardant l'heure
          d'origine à côté.
    \item Noter les hypothèses (dérive constante, fuseau supposé).
\end{enumerate}
\end{boxinfo}

% ============================================================================
\section{Fusionner et Corréler}
\label{sec:c07-correler}
% ============================================================================

Des outils de super-chronologie, comme log2timeline et Plaso, extraient
automatiquement les horodatages de centaines de formats (journaux, registre,
système de fichiers, navigateurs) et les fusionnent dans une seule liste.
Une telle liste compte vite des millions de lignes. La corrélation consiste
à y relier des événements par un élément commun, appelé \termecle{pivot}.

\begin{tableaucours}{L{2.8cm}L{5.0cm}L{4.4cm}}{%
    \tableauHeader{Pivot} &
    \tableauHeader{Exemple à NYANGA} &
    \tableauHeader{Ce qu'il relie}
}
Adresse IP & 203.0.113.47 & Balise de SRV-AD01, exfiltration de SRV-FIC01 \\
Compte & \texttt{adm.kamga} & Sessions sur SRV-CPT01 et SRV-FIC01 depuis
SRV-AD01 \\
Empreinte de fichier & Binaire \texttt{WinUpd\textbackslash svchost.exe} &
Même fichier sur SRV-AD01 et SRV-FIC01 \\
Heure & 14~mars, 1~h~12 & Début de l'exfiltration, échanges volumineux de la
balise \\
Nom & \texttt{cdn-update.example.net} & Requêtes DNS, échanges TLS \\
\end{tableaucours}
\vspace{-0.8em}
\begin{center}{\small\itshape Tableau~7.2. Les pivots de corrélation}\end{center}

Un pivot relie, il ne prouve pas à lui seul un lien de cause. Deux machines
qui parlent à la même adresse peuvent obéir au même attaquant, ou utiliser
le même service légitime. Le lien se confirme par une seconde coïncidence
indépendante~: même adresse \emph{et} même binaire, même compte \emph{et}
même origine.

% ============================================================================
\section{La Chronologie de NYANGA}
\label{sec:c07-nyanga}
% ============================================================================

\begin{boxscenario}{Opération NYANGA. Du correctif au virement}
\begin{tableaucours}{L{2.8cm}L{6.0cm}L{2.2cm}c}{%
    \tableauHeader{Date et heure} &
    \tableauHeader{Événement} &
    \tableauHeader{Source} &
    \tableauHeader{Cote}
}
12/12/2024 & Bulletin de l'éditeur sur VULN-2024-4471 & Éditeur & \OUI{} \\
15/12/2024 & Correctif publié & Éditeur & \OUI{} \\
19/12/2024 & Inscription au catalogue KEV & CISA & \OUI{} \\
02/03/2025, 23~h~48 & Première session \texttt{svc\_backup} depuis
203.0.113.18 & Passerelle (corrigée) & \OUI{} \\
03 au 08/03 & Activité de l'attaquant & Aucune & \NEK{} \\
09/03, 1~h~34 & Session \texttt{svc\_backup} vers SRV-AD01 & Passerelle,
SRV-AD01 & \OUI{} \\
09/03, 1~h~35 & Service \texttt{WinUpdSvc} installé & SRV-AD01 (7045) &
\OUI{} \\
09/03, 1~h~37 & \texttt{lsass.dmp} créé & SRV-AD01 (\$UsnJrnl) & \OUI{} \\
10/03 et suivants & Sessions \texttt{adm.kamga} depuis SRV-AD01 & Serveurs
(4624) & \OUI{} \\
12/03, 22~h~39 à 23~h~06 & Ouverture de \texttt{api.ini}, archive,
suppression & SRV-CPT01 & \OUI{} \\
14/03, 1~h~12 à 3~h~47 & 3,2~Go de SRV-FIC01 vers 203.0.113.47 & Sonde,
NC-07 & \OUI{} \\
14/03, 2~h~05 à 3~h~31 & 23 virements depuis 203.0.113.62 & Interface &
\OUI{} \\
14/03, 7~h~30 & Détection par le service des paiements & Témoignage &
\OUI{} \\
14/03, 10~h~02 & Capture mémoire de SRV-AD01 & Registre des actions &
\OUI{} \\
\end{tableaucours}

\medskip
\textbf{Constatations.} Entre la première session et la détection
s'écoulent 11~jours, 7~heures et 42~minutes. Aucune source ne couvre
l'activité de l'attaquant du 3 au 8~mars~: la passerelle n'enregistre aucune
session de \texttt{svc\_backup}, et aucun serveur du siège ne garde trace de
lui.

\textbf{Analyse} (confiance moyenne). Le trou peut signifier que l'attaquant
a attendu, ou qu'il a agi par un autre chemin que les sources ne voient pas
(une agence, un autre compte). La seconde explication ne peut être écartée.
\end{boxscenario}

\begin{boxdef}
\textbf{\termecle{Temps de présence}} (\textit{dwell time})~: durée entre
la première action établie de l'attaquant dans le système et la détection
de l'incident. C'est une borne inférieure~: une action plus ancienne, non
détectée, l'allongerait.
\end{boxdef}

% ============================================================================
\section{Lire la Chronologie par Tactiques}
\label{sec:c07-tactiques}
% ============================================================================

La chaîne d'intrusion décrite par Lockheed Martin~\cite{killchain2011} et la
matrice ATT\&CK de MITRE~\cite{mitreattack} rangent les actions d'un
attaquant en étapes, ou tactiques. Placer chaque événement de la chronologie
dans une tactique permet de voir ce qui manque.

\begin{tableaucours}{L{3.6cm}L{5.4cm}c}{%
    \tableauHeader{Tactique (ATT\&CK)} &
    \tableauHeader{Événement de NYANGA} &
    \tableauHeader{Cote}
}
Accès initial & Session du 2~mars par la passerelle & \OUI{} \\
Persistance & Service \texttt{WinUpdSvc} & \OUI{} \\
Accès aux identifiants & Copie de la mémoire de \texttt{lsass.exe} & \OUI{} \\
Mouvement latéral & Sessions \texttt{adm.kamga} depuis SRV-AD01 & \OUI{} \\
Collecte & Archive \texttt{cpt.7z}, dossier \texttt{Clients} & \OUI{} \\
Commande et contrôle & Balise de 60~secondes vers 203.0.113.47 & \OUI{} \\
Exfiltration & 3,2~Go le 14~mars & \OUI{} \\
Impact & 23 virements frauduleux & \OUI{} \\
Découverte (reconnaissance interne) & Aucun événement retrouvé & \NEK{} \\
\end{tableaucours}
\vspace{-0.8em}
\begin{center}{\small\itshape Tableau~7.3. La chronologie de NYANGA par tactiques}\end{center}

La dernière ligne est instructive. Un attaquant qui trouve le fichier de clé
d'une application de paiement a forcément exploré le réseau avant. Aucune
trace de cette exploration n'a été trouvée~: elle s'est probablement
déroulée pendant le trou du 3 au 8~mars, avec des moyens que les sources du
siège n'enregistrent pas. La tactique est presque certaine, l'événement
reste \NEK{}.

\begin{boxlecture}{La chronologie au 26 mars}
\textbf{Décompte.} Sur les quatorze lignes de la chronologie, treize sont
cotées \OUI{} et une \NEK{}~; sur les neuf tactiques, huit \OUI{} et une
\NEK{}.

\medskip
\textbf{Lecture.} La chronologie est solide sur ce qu'elle montre, et ses
deux \NEK{} portent sur la même période~: du 3 au 8~mars. C'est là que
l'enquête doit chercher, en élargissant ses sources aux agences et aux
autres comptes, avant de conclure que le temps de présence est de onze
jours.
\end{boxlecture}

% ============================================================================
\section{Ce que la Chronologie demande à chaque Métier}
\label{sec:c07-metiers}
% ============================================================================

\begin{boxmetier}{Analyste en réponse à incident}
\textbf{Construire une chronologie}
\begin{enumerate}
    \item Dresser le tableau des horloges~: fuseau, dérive, précision.
    \item Normaliser en UTC corrigé, en gardant l'heure d'origine.
    \item Fusionner, puis filtrer par pivots.
    \item Confirmer chaque lien par deux coïncidences indépendantes.
    \item Coter chaque ligne et nommer les trous.
\end{enumerate}
\end{boxmetier}

\begin{boxmetier}{Application OPJ/Expert Judiciaire}
Devant un tribunal, une chronologie se présente avec ses incertitudes~: on
écrit « entre 1~h~34 et 1~h~35, à une minute près » plutôt qu'une seconde
précise qu'aucune horloge ne garantit.
\end{boxmetier}

\begin{boxmetier}{RSSI}
Synchroniser toutes les horloges, y compris celles des équipements installés
par des prestataires. La passerelle de NYANGA, seule hors synchronisation,
était aussi le point d'entrée.
\end{boxmetier}

\separateur

\begin{boxresume}
\textbf{Ce qu'il faut retenir de ce chapitre~:}
\begin{itemize}
    \item Avant toute fusion, chaque source reçoit son fuseau, sa dérive et
          sa précision~; la chronologie se construit en UTC corrigé.
    \item Un pivot relie des événements~; un lien de cause demande deux
          coïncidences indépendantes.
    \item Le temps de présence est une borne inférieure.
    \item Ranger la chronologie par tactiques montre ce qui manque.
    \item À NYANGA, le temps de présence établi est de 11~jours, 7~heures
          et 42~minutes~; la période du 3 au 8~mars reste \NEK{}.
\end{itemize}
\end{boxresume}

\begin{boxinfo}
\textbf{Points de contrôle.} Vous avez compris ce chapitre si vous savez~:
\begin{itemize}
    \item corriger une heure d'une dérive mesurée~;
    \item calculer un temps de présence~;
    \item ranger des événements par tactiques et repérer une tactique sans
          événement.
\end{itemize}
\end{boxinfo}

% ============================================================================
\section*{Exercices}
% ============================================================================

\begin{boxexo}[title={\ding{45}\quad Exercice~7.1. Corriger une heure \niveaubase}]
La passerelle avance de 47~secondes. Corrigez et convertissez en UTC les
heures suivantes, lues dans son journal en heure locale~: (a) 02/03/2025
23:48:35~; (b) 09/03/2025 01:35:41~; (c) 14/03/2025 00:58:12.
\end{boxexo}

\begin{boxexo}[title={\ding{45}\quad Exercice~7.2. Le temps de présence \niveaumoyen}]
\begin{enumerate}[(a)]
    \item Vérifiez le calcul du temps de présence de NYANGA.
    \item On découvre dans les journaux d'une agence une connexion de
          \texttt{svc\_backup} le 24~février à 21~h~10. Que devient le temps
          de présence~? Quelle vérification faites-vous d'abord~?
    \item Pourquoi parle-t-on de borne inférieure~?
\end{enumerate}
\end{boxexo}

\begin{boxexo}[title={\ding{45}\quad Exercice~7.3. Un pivot trompeur \niveaumoyen}]
Un poste de la direction a parlé à 203.0.113.47 le 11~mars, à 9~h~00, pendant
quatre secondes.
\begin{enumerate}[(a)]
    \item Ce poste est-il compromis~?
    \item Quelles secondes coïncidences chercheriez-vous~?
    \item Comment cotez-vous la proposition « le poste de la direction est
          compromis » avant et après ces vérifications~?
\end{enumerate}
\end{boxexo}

\begin{boxexo}[title={\ding{45}\quad Exercice~7.4. Combler le trou \niveauexpert}]
Proposez un plan d'investigation pour la période du 3 au 8~mars~: sources à
recueillir, questions, pivots, et cotes attendues selon ce que l'on
trouvera. Indiquez ce qui ferait passer la tactique « découverte » de \NEK{}
à \OUI{}, et ce qui ne permettra jamais de la passer à \NON{}.
\end{boxexo}

\bigskip

\begin{boxinfo}
\textbf{Pour aller plus loin.}
\begin{itemize}
    \item La matrice MITRE ATT\&CK~\cite{mitreattack}.
    \item E.~Hutchins, M.~Cloppert, R.~Amin, sur la chaîne
          d'intrusion~\cite{killchain2011}.
\end{itemize}
\end{boxinfo}

\begin{boxinfo}
\textbf{Transition.} La partie~II a montré comment acquérir et analyser.
Mais qui a le droit de le faire, selon quelles normes, et sous quelles
règles de droit~? C'est l'objet du Chapitre~\ref{chap:normes-droit}.
\end{boxinfo}
